\documentclass[10pt,a4paper]{amsart}
\usepackage{amsmath,amssymb,amsthm,graphicx}
\usepackage{verbatim}
\newenvironment{reference}{\comment}{\endcomment}
%\newenvironment{reference}{}{}
\newenvironment{slogan}{\comment}{\endcomment}
\newenvironment{history}{\comment}{\endcomment}
% For commutative diagrams we use Xy-pic
\usepackage[all]{xy}

% We use 2cell for 2-commutative diagrams.
\xyoption{2cell}
% Theorem environments.
%
\theoremstyle{plain}
\newtheorem{theorem}[subsection]{Theorem}
\newtheorem{proposition}[subsection]{Proposition}
\newtheorem{lemma}[subsection]{Lemma}

\theoremstyle{definition}
\newtheorem{definition}[subsection]{Definition}
\newtheorem{example}[subsection]{Example}
\newtheorem{exercise}[subsection]{Exercise}
\newtheorem{situation}[subsection]{Situation}

\theoremstyle{remark}
\newtheorem{remark}[subsection]{Remark}
\newtheorem{remarks}[subsection]{Remarks}

\numberwithin{equation}{subsection}

% Macros
%
\def\lim{\mathop{\mathrm{lim}}\nolimits}
\def\colim{\mathop{\mathrm{colim}}\nolimits}
\def\Spec{\mathop{\mathrm{Spec}}}
\def\Hom{\mathop{\mathrm{Hom}}\nolimits}
\def\Ext{\mathop{\mathrm{Ext}}\nolimits}
\def\SheafHom{\mathop{\mathcal{H}\!\mathit{om}}\nolimits}
\def\SheafExt{\mathop{\mathcal{E}\!\mathit{xt}}\nolimits}
\def\Sch{\mathit{Sch}}
\def\Mor{\mathop{\mathrm{Mor}}\nolimits}
\def\Ob{\mathop{\mathrm{Ob}}\nolimits}
\def\Sh{\mathop{\mathit{Sh}}\nolimits}
\def\NL{\mathop{N\!L}\nolimits}
\def\CH{\mathop{\mathrm{CH}}\nolimits}
\def\proetale{{pro\text{-}\acute{e}tale}}
\def\etale{{\acute{e}tale}}
\def\QCoh{\mathit{QCoh}}
\def\Ker{\mathop{\mathrm{Ker}}}
\def\Im{\mathop{\mathrm{Im}}}
\def\Coker{\mathop{\mathrm{Coker}}}
\def\Coim{\mathop{\mathrm{Coim}}}

% Boxtimes
%
\DeclareMathSymbol{\boxtimes}{\mathbin}{AMSa}{"02}

%
% Macros for moduli stacks/spaces
%
\def\QCohstack{\mathcal{QC}\!\mathit{oh}}
\def\Cohstack{\mathcal{C}\!\mathit{oh}}
\def\Spacesstack{\mathcal{S}\!\mathit{paces}}
\def\Quotfunctor{\mathrm{Quot}}
\def\Hilbfunctor{\mathrm{Hilb}}
\def\Curvesstack{\mathcal{C}\!\mathit{urves}}
\def\Polarizedstack{\mathcal{P}\!\mathit{olarized}}
\def\Complexesstack{\mathcal{C}\!\mathit{omplexes}}
% \Pic is the operator that assigns to X its picard group, usage \Pic(X)
% \Picardstack_{X/B} denotes the Picard stack of X over B
% \Picardfunctor_{X/B} denotes the Picard functor of X over B
\def\Pic{\mathop{\mathrm{Pic}}\nolimits}
\def\Picardstack{\mathcal{P}\!\mathit{ic}}
\def\Picardfunctor{\mathrm{Pic}}
\def\Deformationcategory{\mathcal{D}\!\mathit{ef}}
\usepackage{cleveref}
\crefname{originalreference}{original reference}{original references}
\Crefname{originalreference}{Original reference}{Original references}
\makeatletter
\newlabel{properties-lemma-Noetherian-topology}{{properties-lemma-Noetherian-topology}{source}{}{sourcepointer-1}{}}
\newlabel{properties-lemma-Noetherian-topology@cref}{{[originalreference][0][]properties-lemma-Noetherian-topology}{[1][0][]source}{}{}{}}
\makeatother
\crefname{originalreference}{original reference}{original references}
\Crefname{originalreference}{Original reference}{Original references}
\makeatletter
\newlabel{topology-lemma-image-Noetherian}{{topology-lemma-image-Noetherian}{source}{}{sourcepointer-2}{}}
\newlabel{topology-lemma-image-Noetherian@cref}{{[originalreference][0][]topology-lemma-image-Noetherian}{[1][0][]source}{}{}{}}
\makeatother
\crefname{originalreference}{original reference}{original references}
\Crefname{originalreference}{Original reference}{Original references}
\makeatletter
\newlabel{topology-lemma-quasi-compact-locally-Noetherian-Noetherian}{{topology-lemma-quasi-compact-locally-Noetherian-Noetherian}{source}{}{sourcepointer-3}{}}
\newlabel{topology-lemma-quasi-compact-locally-Noetherian-Noetherian@cref}{{[originalreference][0][]topology-lemma-quasi-compact-locally-Noetherian-Noetherian}{[1][0][]source}{}{}{}}
\makeatother
\title{The space of a locally Noetherian algebraic stack is locally Noetherian, and Noetherian when the stack is quasi-compact}
\author{The Stacks Project Authors}
\date{stacks-project:stacks-morphisms.tex:local-snapshot-20261012; GNU FDL 1.2 or later}
\begin{document}
\maketitle
\noindent\textit{Editorial note.} All delivered mathematical text is the author's own, copied from the pinned Stacks Project chapter source sources/stacks-algebraic-stacks-morphisms/pinned-original-source.tex; no statement or proof step was rewritten, reordered or removed.\par\medskip
\noindent\textit{Editorial note.} Preamble provenance: the Stacks Project keeps one single shared preamble for every chapter; this card takes it directly from the staged sources/stacks-algebraic-stacks-morphisms/preamble.tex as three verbatim spans (lines 8-12, 14-18 and 154-217), so no derived file was written and nothing was added or redefined. The wrapper class and the base packages are supplied by the assembler.\par\medskip
\noindent\textit{Editorial note.} Bibliography provenance: neither the counted statement nor its proof contains a citation (this lane rescans the whole span before writing the card), so the card carries no bibliography and no entry was reconstructed.\par\medskip
\noindent\textit{Editorial note.} Selection provenance: this card was opened under the file-based control contract (CONTROL.md), with the mechanical pre-card gate run on its statement and proof interval before the card was written.\par\medskip
\noindent\textit{Allowed standard prerequisites (editorial).}\begin{itemize}
\item The statement for schemes that the space of a locally Noetherian scheme is a locally Noetherian topological space. Invoked inside the counted proof (source lines 2266-2268) as a cross-chapter pointer. Provides the property for the presenting scheme.
\item The statement that the image of a locally Noetherian topological space under an open surjective map is locally Noetherian. Invoked inside the counted proof (source lines 2269-2272) as a cross-chapter pointer. Descends local Noetherianity from the presentation to the space of the stack.
\item The statement that a quasi-compact locally Noetherian topological space is Noetherian. Invoked at the end of the counted proof (source lines 2274-2276) as a cross-chapter pointer. Gives the second assertion once quasi-compactness is known.
\end{itemize}\medskip
\ifdefined\hypertarget\hypertarget{sourcepointer-1}{}\fi
\noindent\textit{Source reference.} Original reference properties-lemma-Noetherian-topology: sources/stacks-algebraic-stacks-morphisms/pinned-original-source.tex, source lines 2268--2268. Navigation pointer only: the counted text refers to the Stacks item properties-lemma-Noetherian-topology, which belongs to another chapter of the project and is not part of this card. The pointer carries the target's own label as its number because this chapter's staged snapshot does not contain the other chapter's tag table; its content is not reproduced here.\par\medskip
\ifdefined\hypertarget\hypertarget{sourcepointer-2}{}\fi
\noindent\textit{Source reference.} Original reference topology-lemma-image-Noetherian: sources/stacks-algebraic-stacks-morphisms/pinned-original-source.tex, source lines 2272--2272. Navigation pointer only: the counted text refers to the Stacks item topology-lemma-image-Noetherian, which belongs to another chapter of the project and is not part of this card. The pointer carries the target's own label as its number because this chapter's staged snapshot does not contain the other chapter's tag table; its content is not reproduced here.\par\medskip
\ifdefined\hypertarget\hypertarget{sourcepointer-3}{}\fi
\noindent\textit{Source reference.} Original reference topology-lemma-quasi-compact-locally-Noetherian-Noetherian: sources/stacks-algebraic-stacks-morphisms/pinned-original-source.tex, source lines 2276--2276. Navigation pointer only: the counted text refers to the Stacks item topology-lemma-quasi-compact-locally-Noetherian-Noetherian, which belongs to another chapter of the project and is not part of this card. The pointer carries the target's own label as its number because this chapter's staged snapshot does not contain the other chapter's tag table; its content is not reproduced here.\par\medskip
\begin{lemma}
\label{lemma-Noetherian-topology}
Let $\mathcal{X}$ be an algebraic stack.
\begin{enumerate}
\item If $\mathcal{X}$ is locally Noetherian then $|\mathcal{X}|$
is a locally Noetherian topological space.
\item If $\mathcal{X}$ is quasi-compact and locally Noetherian, then
$|\mathcal{X}|$ is a Noetherian topological space.
\end{enumerate}
\end{lemma}

\begin{proof}
Assume $\mathcal{X}$ is locally Noetherian.
Choose a scheme $U$ and a surjective smooth morphism
$U \to \mathcal{X}$. As $\mathcal{X}$ is locally Noetherian
we see that $U$ is locally Noetherian. By
Properties, Lemma \ref{properties-lemma-Noetherian-topology}
this means that $|U|$ is a locally Noetherian topological space.
Since $|U| \to |\mathcal{X}|$ is open and surjective we conclude that
$|\mathcal{X}|$ is locally Noetherian by
Topology, Lemma \ref{topology-lemma-image-Noetherian}.
This proves (1). If $\mathcal{X}$ is quasi-compact and locally Noetherian,
then $|\mathcal{X}|$ is quasi-compact and locally Noetherian. Hence
$|\mathcal{X}|$ is Noetherian by Topology, Lemma
\ref{topology-lemma-quasi-compact-locally-Noetherian-Noetherian}.
\end{proof}

\end{document}
